\section{Registro persistente de plazos y atención auditada}\label{sec:implementacion-plazos-persistentes}

El seguimiento por expediente compone el catálogo de perfiles, la resolución autorizada de fuentes y el evaluador en una operación persistente. Cada plazo conserva identidad, título, revisión exacta del perfil, selección de insumos, responsable y resultado. Un registro puede quedar bloqueado y carecer de fecha de vencimiento: persistirlo conserva la explicación y permite corregirlo, sin completar datos ausentes ni sustituir el cálculo por una fecha manual. La aplicación de un perfil sigue dependiendo de las declaraciones, condiciones y fuentes descritas en la \cref{sec:implementacion-perfiles-plazos}; su almacenamiento no acredita la interpretación jurídica ni la eficacia del acto inicial.

\subsection{Definición, resultado y evolución de la captura}

El alta y la corrección explícita resuelven nuevamente los insumos y ejecutan el evaluador. El perfil seleccionado debe corresponder a su revisión publicada actual y ser accesible para el expediente. Las fuentes procesales y el calendario conservan, por separado, la revisión elegida y la cabeza observada durante la preparación. Una notificación puede remitir a una revisión de su resolución distinta de la que declara su cabeza; ambas relaciones se conservan. En resultados de audiencia, el acuerdo pertenece a la selección histórica y no se añade artificialmente a la cabeza consultada.

La entrada \texttt{DEVI1} conserva selección, calificación declarada, calendario exacto, cantidad ordenada opcional, declaración de aplicabilidad, incidentes y condiciones en su orden. Se diferencian ausencia, desconocimiento y valores conocidos, incluido un UUID cero. Una cantidad ausente no se rellena con el máximo del perfil. La precisión temporal y el desfase original permanecen explícitos: fecha, minuto y segundo no se intercambian, y la ausencia de desfase no significa UTC.

El formato \texttt{DRES1} conserva el resultado obtenido: requisito, extracción o bloqueo, regla instanciada cuando existe, ancla aritmética, candidata, traza y bloqueos de evaluación. Para días contados mantiene cada clasificación del calendario, su procedencia, explicación y referencias. Para meses conserva el año y mes objetivo y el día solicitado, incluso si no existe homólogo; para horas mantiene la cantidad y los instantes de la operación. Un instante conserva segundos, nanosegundos y desfase. La consulta histórica decodifica esa captura y verifica su vinculación con los insumos; no vuelve a ejecutar la aritmética sobre las cabezas actuales.

El responsable se selecciona entre cuentas activas con permiso de lectura del plazo y acceso al expediente. La captura conserva identificador, correo y rol, sin conceder membresía por la asignación. El estado de atención se representa como pendiente o como actuación declarada con tiempo, relato y localizador. Registrar atención no demuestra presentación válida ni oportuna. Puede conservar tiempo desconocido y no fabrica una hora. La corrección conserva la atención existente; su cambio exige la operación correspondiente.

Cada mutación humana produce una revisión. La corrección exige base y motivo; el cambio de atención y el retiro conservan definición, responsable y cálculo histórico. El retiro es terminal y organizativo, sin declarar nulidad del acto. Las condiciones activo o retirado, atención pendiente o declarada, cálculo bloqueado e instante calculado son dimensiones distintas. Los estados derivados del reloj y la presentación del vencimiento operativo requieren todavía su integración con el servicio y la interfaz.

\subsection{Confirmación transaccional y evidencia histórica}

La migración \texttt{0017} introduce \texttt{case\_deadlines} y \texttt{case\_deadline\_revisions}. La raíz fija el expediente; las revisiones conservan los bytes de entrada y resultado, perfil, dependencias exactas, administración observada, responsable, atención, autor y recibo. Las referencias tipificadas vinculan las revisiones de fuentes, sus cabezas capturadas y el calendario. La proyección SQL de \texttt{DEVI1} recorre estrictamente la representación completa para recuperar esas referencias; no implementa un segundo evaluador.

La confirmación reabre la preparación dentro de una transacción y toma el bloqueo común de auditoría. Revalida identidad activa, permisos, pertenencia, expediente abierto y revisión base. En alta y corrección comprueba además las fuentes vigentes y la elegibilidad actual del responsable; atención y retiro conservan los datos históricos de la captura. El reloj definitivo se consulta después del bloqueo. El estado y el evento de auditoría se confirman juntos; un fallo de validación o escritura revierte la mutación. Los identificadores de operación y la comparación de revisión impiden aceptar un segundo envío como una nueva continuación silenciosa.

Los recibos V1 separan tres huellas. \texttt{DLRV1} vincula el contenido revisado; \texttt{DLST1} añade el contexto administrativo capturado; \texttt{DLTX1} fija actor, expediente, plazo, operación, acción, base, huella revisada y motivo. Esta separación permite que alta o corrección capturen una administración activa más reciente sin alterar el contenido confirmado. Se rechazan retrocesos y cambios de contenido bajo la misma revisión administrativa. Un expediente sin revisiones conserva su contexto original sin atribuirle revisión, fecha o autor ficticios. Atención y retiro mantienen exactamente la administración del cálculo histórico, aunque la administración actual avance.

La autorización del servicio precede al puerto y se repite antes de entregar una consulta o confirmar un envío. PostgreSQL comprueba, a su vez, las condiciones vigentes dentro de la transacción. Estas comprobaciones no convierten Redis y PostgreSQL en una sola transacción. Owner y Litigator pueden gestionar plazos; Paralegal puede consultar, y Client permanece denegado. Litigator y Paralegal requieren pertenencia vigente al expediente. El cierre administrativo conserva las lecturas autorizadas y bloquea las mutaciones humanas.

El lector histórico reconstruye las dependencias capturadas y contrasta sus recibos, huellas y proyecciones. El arranque comprueba esquema, funciones, restricciones, disparadores y permisos; el inventario recorre las revisiones almacenadas. La restauración debe mantener las referencias históricas, los bytes canónicos y la auditoría. Esa obligación se comprueba mediante un recorrido específico y no se deduce de que el esquema pueda crearse en una base vacía.

\subsection{Seguimiento V2 y conservación de las decisiones humanas}

Las migraciones del grupo \texttt{0018} incorporan capturas de seguimiento V2 sin reinterpretar los recibos V1. La nueva captura conserva políticas independientes para perfil, fuente y calendario, estado de revisión, motivos y observaciones verificadas. \texttt{DLTX2} vincula la operación con esas observaciones, ambas huellas del predecesor y una causa explícita. La autoría distingue una cuenta humana del servicio técnico de reevaluación; este último no utiliza una cuenta ficticia ni adopta la identidad de quien modificó la fuente. Un recibo V2 inválido se rechaza y no se vuelve a interpretar como V1.

El manifiesto \texttt{DLOB1} identifica las cabezas examinadas de perfil, fuente, calendario y padre de notificación cuando corresponde. Cada observación compromete revisión, ámbito, recibo de origen y evidencia histórica completa. La cabeza del padre permanece separada de la resolución histórica vinculada a la notificación. El lector reconstruye esas revisiones exactas y comprueba también la administración observada; una coincidencia de huellas aisladas no sustituye esta verificación. Una revisión administrativa igual debe conservar todo su contenido, incluidos autor y componentes temporales, y una revisión posterior no puede retroceder.

Las políticas permiten conservar una selección fija, seguir sus cambios o expresar que la decisión sigue sin determinarse. Un cambio de fuente o perfil seguido exige revisión humana y mantiene el cálculo histórico. El retiro de una dependencia o una política no declarada conserva motivos explícitos. El consumidor combina los motivos con los ya pendientes, sin borrarlos ni aceptar por el usuario. Sólo el cambio de un calendario seguido y publicado puede recalcular una captura previamente aceptada cuando la revisión resultante también permanece aceptada; conserva para ello el perfil, la fuente y las calificaciones históricas. Si un cambio simultáneo exige revisión, conserva el cálculo anterior.

El vencimiento operativo se obtiene únicamente de un plazo activo con revisión aceptada y cálculo que ya contiene un instante final. Una captura pendiente, una captura heredada sin política declarada o un plazo retirado carece de vencimiento operativo. La existencia de fecha en V1 no expresa aceptación. La conciliación del legado incorpora seguimiento sin inventar esa decisión y conserva cálculo, responsable y atención. El almacenamiento permite capturas humanas V2, pero el servicio humano utilizado por HTTP y Qadra sigue produciendo V1.

\subsection{Consumo durable y confirmación de resultados}

El adaptador consumidor procesa los trabajos persistidos por el despachador descrito en la \cref{sec:implementacion-perfiles-plazos}. Cada ejecución toma el bloqueo común de auditoría mediante una transacción \texttt{READ COMMITTED}, selecciona un trabajo elegible y autentica su operación, causa y revisión de origen. Después carga la cabeza verificada del plazo. Esta secuencia conserva cualquier corrección humana confirmada previamente. El cierre del expediente y la inactividad del responsable no excluyen el seguimiento técnico; tampoco conceden al consumidor permisos para registrar decisiones humanas.

Las migraciones \texttt{0020} agregan resultados e intentos inmutables. Cuando procede una revisión técnica, revisión, resultado y auditoría se confirman en la misma transacción. El resultado referencia las dos huellas de su base exacta y las de la revisión producida. Las restricciones impiden confirmar un resultado de revisión sin su revisión correspondiente y, mediante una comprobación diferida, una revisión técnica sin resultado. La operación reservada por el trabajo no puede reutilizarse desde el puerto humano.

La ejecución también puede concluir sin crear una revisión. Sus motivos distinguen plazo retirado, legado ya inicializado, dependencia que dejó de estar seleccionada y revisión de fuente ya observada. Los dos primeros se justifican mediante base y causa verificadas; los otros conservan además las observaciones examinadas y el compromiso de administración completa. La consulta reconstruye esa evidencia histórica, sin sustituirla por las cabezas actuales. Los resultados siguen siendo consultables después de una corrección posterior, y un trabajo ya confirmado no produce una segunda revisión al consultarlo de nuevo.

\subsection{Intentos fallidos y recuperación}

Una ejecución fallida revierte su transacción antes de intentar registrar el fallo en otra transacción auditada. Cada intento conserva identidad propia, número consecutivo dentro del trabajo, categoría, código estable, base comprobada cuando existe y tiempos de fallo y próximo reintento. La base sólo se declara comprobada después de verificar su recibo. Los errores tipados diferencian indisponibilidad, espera de bloqueo e interrupción de transacción frente a evidencia almacenada o trabajo inconsistentes. Cuando el puerto no conserva información tipada suficiente, se usa un código genérico; no se deduce la categoría del texto de un mensaje SQL.

Los fallos transitorios esperan uno, dos, cuatro y sucesivos segundos hasta un máximo de 300 para la misma base. Las inconsistencias esperan 3600 segundos. Una cabeza posterior deja sin efecto la espera asociada a la base anterior, de modo que el consumidor examina la nueva decisión humana. Otros trabajos elegibles pueden continuar. El intento y su auditoría se confirman juntos; si se pierde la respuesta, se consulta su identidad exacta y el resultado del trabajo antes de informar una confirmación. La imposibilidad de registrar o conciliar se comunica como error, sin simular un intento durable.

La ausencia de trabajo elegible no implica que todos los plazos estén actualizados, pues puede haber intentos en espera. La apertura y la reconexión verifican catálogo, permisos e inventario completo, incluidos todos los intentos y las revisiones técnicas vinculadas. Una inconsistencia persistente impide abrir hasta repararla. El respaldo conserva conjuntamente trabajos, resultados, intentos, plazos, fuentes y auditoría; la restauración exige comprobar los recibos y continuar desde el progreso conservado. Los límites de un segundo para espera de bloqueo y cinco segundos por sentencia no acotan la duración completa de apertura, evaluación o ejecución. El contrato y sus límites se documentan en \rutarepo{docs/deadline-worker.md} y la decisión en \rutarepo{docs/adr/0037-durable-deadline-reevaluation.md}.

\subsection{Consulta HTTP y alcance operativo}

La API organiza el recurso bajo la colección de plazos de cada expediente. Ofrece preparación, alta, corrección, atención, retiro, listado, detalle, revisión exacta e historia. Las mutaciones confirman la huella del envío preparado. Las páginas son acotadas y el historial devuelve metadatos y recibos ligeros; la revisión exacta proporciona el cálculo completo. Así se evita repetir perfiles y trazas voluminosos en cada entrada de una página de historia.

La proyección JSON conserva la entrada y el resultado estructurados, incluidos bloqueos, trazas, precisión y fuentes, sin recurrir a representaciones de depuración ni recalcularlos para presentarlos. El recorrido HTTP integrado comprueba las operaciones V1 sobre servicios reales y conserva las respuestas exactas después de restaurar la base. Esta evidencia se registra separadamente de las pruebas del consumidor local. La interfaz Qadra permite operar el registro por expediente. El despachador y el consumidor aún no se componen en \texttt{serve}; la evolución del servicio humano, HTTP y Qadra para exponer políticas, revisión pendiente, causa y autoría V2 conserva su integración pendiente. Tampoco se han incorporado activación automática, agenda combinada ni alertas internas y por correo.


\subsection{Interfaz por expediente y selección autorizada}

Qadra incorpora una vista de plazos para el personal autorizado dentro del expediente seleccionado. La colección pagina los registros activos o retirados; cada detalle distingue el vencimiento calculado, los bloqueos y la atención declarada. El historial recupera una revisión exacta con su definición, resultado, autor y recibo. Su consulta no sustituye la captura histórica por la evaluación de las fuentes actuales. La presentación conserva nanosegundos y desfase del instante, y muestra por separado el calendario y la fuente seleccionados y las cabezas observadas al registrar.

El formulario utiliza selectores de perfiles, fuentes, calendarios y responsables. El catálogo de perfiles puede ser global o privado del expediente; el operador consulta definición, condiciones, referencias y ejemplos antes de confirmar su selección. Las fuentes comprenden resoluciones, notificaciones con padre histórico y resultados de audiencia con acuerdo explícito cuando corresponde. La declaración de desconocimiento exige un motivo. Los valores de aplicabilidad, incidentes y condiciones no se completan como verdaderos o falsos de manera implícita. Una cantidad concedida ausente permanece ausente, aunque el perfil establezca un máximo.

El selector de responsables expone únicamente identificador, correo y rol de cuentas elegibles. Incluye Owner activos con acceso global y Litigator o Paralegal activos asignados al expediente; no añade Client ni expone credenciales. La consulta aplica autorización y filtros antes de paginar, confirma su auditoría y admite lectura en expedientes cerrados. La selección no reserva permisos: el alta o la corrección vuelve a comprobar la elegibilidad durante la confirmación transaccional.

La preparación permite revisar el cálculo, la candidata y los motivos de bloqueo antes del envío. El usuario confirma de manera explícita un resultado bloqueado; registrar esa situación no le atribuye fecha final. La corrección exige una base vigente y conserva el borrador ante un conflicto. El operador consulta la versión actual y decide si adopta esa base antes de preparar de nuevo. Si la respuesta del envío se pierde, la interfaz consulta su revisión exacta y contrasta la operación y las huellas; no repite automáticamente una mutación cuyo resultado aún desconoce.

El cierre administrativo deshabilita las modificaciones y mantiene las lecturas autorizadas. Un cambio de expediente o de sesión descarta las respuestas tardías de la vista anterior. La pérdida de autorización elimina sus datos visibles. Las tablas de trazas preservan el orden y el fundamento de cada paso; en pantallas estrechas mantienen columnas legibles mediante desplazamiento horizontal dentro de su contenedor. Estos mecanismos de la interfaz V1 permiten revisar la captura conservada; la presentación del seguimiento V2, los estados del reloj y los avisos requieren sus propios recorridos de integración.
